\documentclass{tex-article}
\addbibresource{bibliography/refs.bib}
\title{TAReL: Temporal Adversarial Reconstruction of High-Frequency
       Latent Spaces in Financial Markets}
\author[1]{J. Francisco Munoz-Elguezabal}
\affil[1]{ITESO --- The Jesuit University of Guadalajara}
\date{}

\begin{document}\maketitle

\begin{abstract}
We introduce \textbf{TAReL}, a pipeline that studies cryptocurrency
limit-order-book microstructure through the geometry of recurrence
plots and the seed-to-seed reproducibility of generative models
trained on those plots. Stage~1 reconstructs top-25-level state
vectors at 250\,ms resolution under three parallel clocks (physical,
event, information time), buckets them subject to a 60-minute
parquet-boundary constraint, and writes per-bucket binary recurrence
plots calibrated to a locked $RR^{*}{=}5\%$ target. Stage~2 trains a
$128$-dimensional W-GAN-GP on the resulting $64{\times}64$ tensors
and streams four training diagnostics into a locked four-axis
Stability Score. We report run \texttt{00000000c002daaaX} on two
cells, \texttt{binance:BTC-USDC} and \texttt{kraken:BTC-USDC},
across all three clocks. The Stage-1 RQA characterisation produces
distinct $DET$ and $LAM$ distributions across clocks at a fixed
cell, supporting the conjecture that recurrence geometry is a
clock-dependent fingerprint of liquidity dynamics. Stage~2 trained
four (cell, clock) pairs to the first checkpoint; we report EM-gap
convergence and a 2-axis partial Stability Score restricted to the
axes the present training horizon populates. A third cell
(\texttt{coinbase:BTC-USDC}) is excluded due to a stage-one
RR-bisection defect tracked separately.
\end{abstract}

\input{content/sections-v2/00-introduction/00-introduction}
\input{content/sections-v2/01-literature-review/00-paper-review}
\input{content/sections-v2/01-literature-review/01-paper-review}
\input{content/sections-v2/02-problem-solution/00-problem}
\input{content/sections-v2/02-problem-solution/01-solution}
\input{content/sections-v2/03-data-methods/00-data}
\input{content/sections-v2/03-data-methods/01-methods}
\input{content/sections-v2/04-results-discussion/00-results}
\input{content/sections-v2/04-results-discussion/01-discussion}
\printbibliography

\end{document}
